% Options for packages loaded elsewhere
\PassOptionsToPackage{unicode}{hyperref}
\PassOptionsToPackage{hyphens}{url}
%
\documentclass[
  a4paper,
]{scrbook}

\usepackage{amsmath,amssymb}
\usepackage{setspace}
\usepackage{iftex}
\ifPDFTeX
  \usepackage[T1]{fontenc}
  \usepackage[utf8]{inputenc}
  \usepackage{textcomp} % provide euro and other symbols
\else % if luatex or xetex
  \usepackage{unicode-math}
  \defaultfontfeatures{Scale=MatchLowercase}
  \defaultfontfeatures[\rmfamily]{Ligatures=TeX,Scale=1}
\fi
\usepackage{lmodern}
\ifPDFTeX\else  
    % xetex/luatex font selection
\fi
% Use upquote if available, for straight quotes in verbatim environments
\IfFileExists{upquote.sty}{\usepackage{upquote}}{}
\IfFileExists{microtype.sty}{% use microtype if available
  \usepackage[]{microtype}
  \UseMicrotypeSet[protrusion]{basicmath} % disable protrusion for tt fonts
}{}
\makeatletter
\@ifundefined{KOMAClassName}{% if non-KOMA class
  \IfFileExists{parskip.sty}{%
    \usepackage{parskip}
  }{% else
    \setlength{\parindent}{0pt}
    \setlength{\parskip}{6pt plus 2pt minus 1pt}}
}{% if KOMA class
  \KOMAoptions{parskip=half}}
\makeatother
\usepackage{xcolor}
\usepackage[left=2.54cm,right=2.54cm]{geometry}
\setlength{\emergencystretch}{3em} % prevent overfull lines
\setcounter{secnumdepth}{5}
% Make \paragraph and \subparagraph free-standing
\makeatletter
\ifx\paragraph\undefined\else
  \let\oldparagraph\paragraph
  \renewcommand{\paragraph}{
    \@ifstar
      \xxxParagraphStar
      \xxxParagraphNoStar
  }
  \newcommand{\xxxParagraphStar}[1]{\oldparagraph*{#1}\mbox{}}
  \newcommand{\xxxParagraphNoStar}[1]{\oldparagraph{#1}\mbox{}}
\fi
\ifx\subparagraph\undefined\else
  \let\oldsubparagraph\subparagraph
  \renewcommand{\subparagraph}{
    \@ifstar
      \xxxSubParagraphStar
      \xxxSubParagraphNoStar
  }
  \newcommand{\xxxSubParagraphStar}[1]{\oldsubparagraph*{#1}\mbox{}}
  \newcommand{\xxxSubParagraphNoStar}[1]{\oldsubparagraph{#1}\mbox{}}
\fi
\makeatother


\providecommand{\tightlist}{%
  \setlength{\itemsep}{0pt}\setlength{\parskip}{0pt}}\usepackage{longtable,booktabs,array}
\usepackage{calc} % for calculating minipage widths
% Correct order of tables after \paragraph or \subparagraph
\usepackage{etoolbox}
\makeatletter
\patchcmd\longtable{\par}{\if@noskipsec\mbox{}\fi\par}{}{}
\makeatother
% Allow footnotes in longtable head/foot
\IfFileExists{footnotehyper.sty}{\usepackage{footnotehyper}}{\usepackage{footnote}}
\makesavenoteenv{longtable}
\usepackage{graphicx}
\makeatletter
\newsavebox\pandoc@box
\newcommand*\pandocbounded[1]{% scales image to fit in text height/width
  \sbox\pandoc@box{#1}%
  \Gscale@div\@tempa{\textheight}{\dimexpr\ht\pandoc@box+\dp\pandoc@box\relax}%
  \Gscale@div\@tempb{\linewidth}{\wd\pandoc@box}%
  \ifdim\@tempb\p@<\@tempa\p@\let\@tempa\@tempb\fi% select the smaller of both
  \ifdim\@tempa\p@<\p@\scalebox{\@tempa}{\usebox\pandoc@box}%
  \else\usebox{\pandoc@box}%
  \fi%
}
% Set default figure placement to htbp
\def\fps@figure{htbp}
\makeatother
% definitions for citeproc citations
\NewDocumentCommand\citeproctext{}{}
\NewDocumentCommand\citeproc{mm}{%
  \begingroup\def\citeproctext{#2}\cite{#1}\endgroup}
\makeatletter
 % allow citations to break across lines
 \let\@cite@ofmt\@firstofone
 % avoid brackets around text for \cite:
 \def\@biblabel#1{}
 \def\@cite#1#2{{#1\if@tempswa , #2\fi}}
\makeatother
\newlength{\cslhangindent}
\setlength{\cslhangindent}{1.5em}
\newlength{\csllabelwidth}
\setlength{\csllabelwidth}{3em}
\newenvironment{CSLReferences}[2] % #1 hanging-indent, #2 entry-spacing
 {\begin{list}{}{%
  \setlength{\itemindent}{0pt}
  \setlength{\leftmargin}{0pt}
  \setlength{\parsep}{0pt}
  % turn on hanging indent if param 1 is 1
  \ifodd #1
   \setlength{\leftmargin}{\cslhangindent}
   \setlength{\itemindent}{-1\cslhangindent}
  \fi
  % set entry spacing
  \setlength{\itemsep}{#2\baselineskip}}}
 {\end{list}}
\usepackage{calc}
\newcommand{\CSLBlock}[1]{\hfill\break\parbox[t]{\linewidth}{\strut\ignorespaces#1\strut}}
\newcommand{\CSLLeftMargin}[1]{\parbox[t]{\csllabelwidth}{\strut#1\strut}}
\newcommand{\CSLRightInline}[1]{\parbox[t]{\linewidth - \csllabelwidth}{\strut#1\strut}}
\newcommand{\CSLIndent}[1]{\hspace{\cslhangindent}#1}

% Palatino font:
\usepackage[T1]{fontenc}
\usepackage[utf8]{inputenc}
\usepackage{palatino}

\usepackage{mathpazo}

\usepackage{upquote}

% front page
\usepackage[colorlinks]{hyperref} % for text color on title
\usepackage{titling} % makes \thetitle etc available...
\usepackage{lipsum}
\usepackage{tikz}
\tikzstyle{bag} = [align=center] % allows new lines in text below
\AtBeginDocument{\thispagestyle{empty}
\definecolor{titlecolor}{HTML}{1C2404}
\begin{tikzpicture}[remember picture,overlay]
\draw (current page.center) node[inner sep=0] {\resizebox{!}{350mm}{\includegraphics{highres_banner.pdf}}};
\draw (current page.center) node[inner sep=0, opacity=0.5] {\resizebox{!}{100mm}{\includegraphics{ausegl_sort.pdf}}};
% \draw (current page.center) node[bag] {  
        \draw (17, 1.7) node[bag, align=right, anchor=north east] {  
    % \mytitlefont 
    \color{titlecolor} \Large \textbf{Thesis} \\ 
    \\ 
    % \mytitlefont 
    \color{titlecolor} \Huge \textbf{\thetitle} \\ 
    \\ 
    % \mytitlefont
    \color{titlecolor} \huge \theauthor \\
    \\
    % \mytitlefont 
    \color{titlecolor} \Large \thedate 
    };
\end{tikzpicture}
\clearpage}

% back side
\AtEndDocument{\clearpage\thispagestyle{empty}
\begin{tikzpicture}[remember picture,overlay]
\draw (current page.center) node[inner sep=0] {\resizebox{!}{350mm}{\includegraphics{highres_banner.pdf}}};
% \draw (current page.center) node[bag] {  
    \draw (17, 1.7) node[bag, align=right, anchor=north east] {  
        % \mytitlefont 
        \color{titlecolor} \large Bioinformatics Research Centre \\ 
        % \mytitlefont 
        \color{titlecolor} \large Department of Molecular Biology and Genetics \\ 
        % \mytitlefont 
        \color{titlecolor} \large Aarhus University \\
        % \mytitlefont 
        \color{titlecolor} \large Universitetsbyen 81 \\
        % \mytitlefont 
        \color{titlecolor} \large 8000 Aarhus C \\
        % \mytitlefont 
        \color{titlecolor} \large Denmark
        };

\end{tikzpicture}}


% % picture (AU logo in title page)
% \usepackage{titlepic}
% % \usepackage{graphicx}
% \titlepic{\resizebox{!}{50mm}{\includegraphics[width=\textwidth]{ausegl_sort.pdf}}}



% \pagestyle{plain} % no running header

\usepackage[many]{tcolorbox}
\definecolor{quotegray}{HTML}{505050}
\newtcolorbox{myquote}{%
    enhanced jigsaw, 
    breakable,      % allow page breaks
    frame hidden,   % hide the default frame
    left=1cm,       % left margin
    right=1cm,      % right margin
    colback=white,
    fontupper=\color{quotegray},
    overlay={%
        \node [scale=4,
            text=lightgray,
            inner sep=0pt,] at ([xshift=0.5cm,yshift=-0.7cm]frame.north west){``}; 
        \node [scale=4,
            text=lightgray,
            inner sep=0pt,] at ([xshift=-0.5cm]frame.south east){''};  
            },
        % paragraph skips obeyed within tcolorbox
                parbox=false,
}
% redefine the 'quote' environment to use this 'myquote' environment
\renewenvironment{quote}{\begin{myquote}}{\end{myquote}}

% add a dummy abstract environment that is defined in quarto manuscript 
% but not in the quarto book that we run first to give it a slot in the 
% side bar
\ifcsmacro{abstract}{}{
  \let\endmyenvironment\undefined%
  \newenvironment{abstract}{\chapter*{Abstract}}{}
}

% \usepackage{framed} % not sure i need this anymore

% \usepackage[T1]{fontenc}
% \usepackage{inconsolata}

% % I have to only define Shaded if it is already defined.
% % The reason is that pandoc does not define the macro if there are not code blocks in the a markdown file.
% \ifx \@Shaded \@empty

% \renewcommand{\KeywordTok}[1]{\textcolor[rgb]{0, 0, 0}{\textbf{{#1}}}} % def and or not reg
% \renewcommand{\BuiltInTok}[1]{\textcolor[rgb]{0.373, 0.298, 0.580}{\textbf{{#1}}}} % print open 
% \renewcommand{\VariableTok}[1]{\textcolor[rgb]{0.141, 0.392, 0.824}{\textbf{{#1}}}}
% \renewcommand{\OperatorTok}[1]{\textcolor[rgb]{0, 0, 0}{{#1}}} % def and or not reg
% \renewcommand{\DataTypeTok}[1]{\textcolor[rgb]{1.0,0.13,0.00}{{#1}}}
% \renewcommand{\DecValTok}[1]{\textcolor[rgb]{0.655, 0.498, 0.161}{{#1}}}
% \renewcommand{\BaseNTok}[1]{\textcolor[rgb]{0.259, 0.592, 0.596}{{#1}}}
% \renewcommand{\FloatTok}[1]{\textcolor[rgb]{0.655, 0.498, 0.161}{{#1}}}
% \renewcommand{\CharTok}[1]{\textcolor[rgb]{0.678,0.141,0.098}{{#1}}}
% \renewcommand{\StringTok}[1]{\textcolor[rgb]{0.678,0.141,0.098}{{#1}}}
% \renewcommand{\CommentTok}[1]{\textcolor[rgb]{0.135, 0.134, 0.133}{{#1}}}
% \renewcommand{\OtherTok}[1]{\textcolor[rgb]{0.00,0.44,0.13}{{#1}}}
% \renewcommand{\AlertTok}[1]{\textcolor[rgb]{1.00,0.00,0.00}{\textbf{{#1}}}}
% \renewcommand{\FunctionTok}[1]{\textcolor[rgb]{0.549, 0.102, 0.063}{\textbf{{#1}}}}  % function name
% \renewcommand{\RegionMarkerTok}[1]{{#1}}
% \renewcommand{\ErrorTok}[1]{\textcolor[rgb]{1.00,0.00,0.00}{\textbf{{#1}}}}
% \renewcommand{\NormalTok}[1]{\textcolor[rgb]{0, 0, 0}{{#1}}}


% \else
%   % no code blocks with markup...
% \fi



\usepackage{etoolbox}
\makeatletter
\g@addto@macro{\appendix}{%
  \patchcmd{\@@makechapterhead}% <cmd>
    {\endgraf\nobreak\vskip.5\baselineskip}% <search>
    {\hspace*{-.5em}:\space}% <replace>
    {}{}% <success><failure>
  \patchcmd{\@chapter}% <cmd>
    {\addchaptertocentry{\thechapter}}% <search>
    {\addchaptertocentry{Appendix~\thechapter:}}% <replace>
    {}{}% <success><failure>
  \addtocontents{toc}{%
    \protect\patchcmd{\protect\l@chapter}% <cmd>
      {1.5em}% <search>
      {6.5em}% <replace>
      {}{}}% <success><failure>
}
\renewcommand{\autodot}{}% Remove all end-of-counter dots
\makeatother


% restart chapter numbers in each part
\makeatletter
\@addtoreset{chapter}{part}
\makeatother


% \usepackage{chngcntr}
% \counterwithin*{subsubsection}{chapter}
% \counterwithout*{subsubsection}{section}
% \counterwithout*{subsubsection}{subsection}

% %  % KMT only use subsubsection number (this increments though the book 
% %  % when we supress numbering with  {.unnumbered} after each header except Exerisices
% \renewcommand\thesubsubsection{\arabic{section}}
% \renewcommand\thesubsubsection{\arabic{subsection}}
% \renewcommand\thesubsubsection{\arabic{chapter}-\arabic{subsubsection}}


\renewcommand*{\chapterformat}{%
\textcolor[rgb]{0.7, 0.7, 0.7}{\thechapter}\autodot\enskip%
}
\renewcommand*{\sectionformat}{%
\textcolor[rgb]{0.7, 0.7, 0.7}{\thesection}\autodot\enskip%
}
\renewcommand*{\subsectionformat}{%
\textcolor[rgb]{0.7, 0.7, 0.7}{\thesubsection}\autodot\enskip%
}
\renewcommand*{\subsubsectionformat}{%
\textcolor[rgb]{0.7, 0.7, 0.7}{\thesubsubsection}\autodot\enskip%
}



% % prevent latex from "floating" the figures
% \usepackage{float}
% \let\origfigure\figure
% \let\endorigfigure\endfigure
% \renewenvironment{figure}[1][2] {
%     \expandafter\origfigure\expandafter[H]
% } {
%     \endorigfigure
% }

\usepackage{xcolor}
% \let\oldtextit\textit 
% \renewcommand\textit[1]{\oldtextit{\color{blue}#1}}
\let\oldemph\emph
\renewcommand\emph[1]{\oldemph{\color{gray}#1}}

\makeatletter
\@ifpackageloaded{caption}{}{\usepackage{caption}}
\AtBeginDocument{%
\ifdefined\contentsname
  \renewcommand*\contentsname{Table of contents}
\else
  \newcommand\contentsname{Table of contents}
\fi
\ifdefined\listfigurename
  \renewcommand*\listfigurename{List of Figures}
\else
  \newcommand\listfigurename{List of Figures}
\fi
\ifdefined\listtablename
  \renewcommand*\listtablename{List of Tables}
\else
  \newcommand\listtablename{List of Tables}
\fi
\ifdefined\figurename
  \renewcommand*\figurename{Figure}
\else
  \newcommand\figurename{Figure}
\fi
\ifdefined\tablename
  \renewcommand*\tablename{Table}
\else
  \newcommand\tablename{Table}
\fi
}
\@ifpackageloaded{float}{}{\usepackage{float}}
\floatstyle{ruled}
\@ifundefined{c@chapter}{\newfloat{codelisting}{h}{lop}}{\newfloat{codelisting}{h}{lop}[chapter]}
\floatname{codelisting}{Listing}
\newcommand*\listoflistings{\listof{codelisting}{List of Listings}}
\makeatother
\makeatletter
\makeatother
\makeatletter
\@ifpackageloaded{caption}{}{\usepackage{caption}}
\@ifpackageloaded{subcaption}{}{\usepackage{subcaption}}
\makeatother

\usepackage{bookmark}

\IfFileExists{xurl.sty}{\usepackage{xurl}}{} % add URL line breaks if available
\urlstyle{same} % disable monospaced font for URLs
\hypersetup{
  pdftitle={Chromatin Compartments and Selection on X},
  pdfauthor={Søren Jørgensen},
  pdfkeywords={Hi-C, Chromatin Compartments, Selection},
  hidelinks,
  pdfcreator={LaTeX via pandoc}}


\title{Chromatin Compartments and Selection on X}
\usepackage{etoolbox}
\makeatletter
\providecommand{\subtitle}[1]{% add subtitle to \maketitle
  \apptocmd{\@title}{\par {\large #1 \par}}{}{}
}
\makeatother
\subtitle{How Edges of Active Chromatin Align with Selection Regions in
Primates}
\author{Søren Jørgensen}
\date{2024-11-26}

\begin{document}
\frontmatter
\cleardoublepage
\thispagestyle{empty}
{\centering
\hbox{}\vskip 0cm plus 1fill

{%\mytitlefont 
\Huge\bfseries Chromatin Compartments and Selection on X \par}
\vspace{3ex}
{\Large\bfseries How Edges of Active Chromatin Align with Selection
Regions in Primates \par}
\vspace{6ex}

    {\Large\bfseries Søren Jørgensen \par}
    %
%
\vskip 0cm plus 2fill

{\resizebox{!}{50mm}{\includegraphics[width=\textwidth]{ausegl_sort.pdf}} \par}
\vskip 0cm plus 2fill

{\bfseries\Large\textit{} \par}
\vspace{3ex}

{\bfseries\large 2024-11-26 \par}
\vspace{3ex}

\vspace{12ex}
{\small Submitted in fulfillment of the requirements
of the degree of  \par}
\pagebreak

\begin{quote}
\raggedright    
3-dimensional structure of chromatin brings light onto the mystery of
selfish genes. \par
\raggedleft    
    {- Søren Jørgensen \par}
%    
\end{quote}
}

\renewcommand*\contentsname{Table of contents}
{
\setcounter{tocdepth}{1}
\tableofcontents
}

\setstretch{1}
\mainmatter
\newpage{}

\chapter{Introduction}\label{introduction}

\section{Sexual reproduction (spermatogenesis,
meiosis)}\label{sexual-reproduction-spermatogenesis-meiosis}

The production of gametes in a sexually reproducing organism is a highly
complex process that involves numeruous elements. Spermatogenesis, the
process of forming male gametes, involves four stages of differentiation
from a germ cell through \emph{spermatogonia}, \emph{pachytene
spermatocyte}, and \emph{round spermatids} to \emph{spermatozoa}, or
\emph{sperm} (Wang et al. 2019), and it is the very basis of male
reproduction. The specialized cell division of meiosis neatly handles
the pairing, recombination, and segregation of homologous chromosomes,
thereby ensuring proper genetic distribution. Deeply understanding the
steps of molecular steps of reproduction and how our genetic material is
inherited is essential in biology, bringing insight to areas such as
speciation, population diversity, and (male) infertility.

\section{Selfish genes (and
randomness)}\label{selfish-genes-and-randomness}

The conventional story of meiosis in gametogenesis is one of random
segregation of the sex chromosomes. They split into haploid gametes,
where each chromosome has an equal chance of being passed on to a
gamete. That seems like a fair game, but what if some genes are cheating
the system by making others less viable. A meiotic driver is a selfish
gene element that modulates meiosis and preferentially transmits its own
allele through meiosis, regardless of the downstream fitness effects it
may have (good or bad) on the organism it is part of. This phenomenon
challenges the traditional understanding of selection, extending its
scope beyond the fitness effects on an organism to include selective
pressures at the molecular level. For example, if some genes on the X
chromosome create a disadvantage for gametes that \emph{do not} contain
those genes, making sure the Y chromosome is not as viable as the X,
resulting in a sex imbalance and possibly numerous other downstream
effects. That is exactly what is coined \emph{sex chromosome meiotic
drive} (Jaenike 2001), a result of selfish genetic elements.

Motivated by previous results in the Munch Research group (Munch 2024)
on hybrid incompatibility and extended common haplotypes (Skov et al.
2023; Sørensen et al. 2023) that could be explained by meiotic drive, we
wanted to investigate how these patterns correlate with chromatin
compartments.

\section{High-Throughput Chromosome Conformation Capture
(Hi-C)}\label{high-throughput-chromosome-conformation-capture-hi-c}

Our DNA can be divided into different orders of structure. \emph{3C}
focus on identifying the highest orders of organization inside the
nucleus, that is, when the 30 nm thick coil of chromatin fibers folds
into loops, Topologically Associating Domains (TADs), and chromatin
compartments. Here, we narrow our focus on the largest of the
structures, \emph{compartments}, that is known to determine availability
to transcription factors, thus making an \emph{A} compartment
\emph{active}---and the \emph{B} compartment \emph{inactive}. The
introduction of the Hi-C (high-throughput 3C) method (Lieberman-Aiden et
al. 2009) opened new possibilities for exploring the three-dimensional
organization of the genome.

\newpage{}

\chapter{Methods}\label{methods}

In this project, we formulate two objectives:

\textbf{A}: Reproduce the Hi-C interaction maps and eigendecomposition
from (Wang et al. 2019), with some modifications. We briefly use
\emph{HiCExplorer}, but change the analyses to use the \emph{Open2C
Ecosystem} (Open Chromosome Collective 2024) which have a Pyton API as
well as command-line functions, which can be paired very well with
Jupyter Notebooks. The majority of the data analysis was run with a
\emph{gwf} workflow, and the commands that were visually inspected were
run in Jupyter Notebooks.

\textbf{B} Compare with regions of selection that are found in
\emph{papio anubis}, and maybe in \emph{human} too. Investigate the
biological meaning of the results.

All computations were performed on GenomeDK (GDK) {[}ref{]}, an HPC
cluster located on Aarhus Uninversity, and most of the processing of the
data was nested into a \emph{gwf} workflow {[}ref{]}, a workflow manager
developed at GDK. I would like to thank GDK and Aarhus University for
providing computational resources and support that contributed to these
research results.

The whole of this project is carried out with reproducibility in mind,
so an effort (and quite a significant amount of time) has been put into
documenting code and organizing the project for readbility and
transparency through a Quarto project {[}ref{]}. Therefore, all code,
virtual environments and text is made available as a Quarto book,
rendered directly from the GitHub repository with GitHub Pages {[}{]}.
To make this possible, the Quarto documentation has been extensively
studied and discussed with \emph{KMT} {[}ref, aknowledge{]}.

\section{Downloading Data and Project
Structure}\label{downloading-data-and-project-structure}

To reproduce the results from (Wang et al. 2019), I chose to use their
raw data directly from the SRA portal {[}ref{]}. I filtered the data to
contain all their paired-end Hi-C reads, and included only macaque
samples. The data set also contains RNAseq data, and the same tissues
for both macaque and mouse. The meta data for the data set was extracted
into a runtable \texttt{SRA-runtable.tsv}. To get an overview of the
data accessions used in this analysis, we will first summarize the
runtable that contains the accession numbers and some metadata for each
sample (Table~\ref{tbl-runtable-summary}). It adds up to
\textasciitilde1Tb of compressed \texttt{fastq} files, holding
\textasciitilde9.5 billion reads, roughly evenly spread on the 5 tissue
types.

\begin{longtable}[]{@{}lllll@{}}

\caption{\label{tbl-runtable-summary}Summary of the data accessions used
in this analysis}

\tabularnewline

\caption{}\label{T_55add}\tabularnewline
\toprule\noalign{}
~ & source\_name & GB & Bases & Reads \\
\midrule\noalign{}
\endfirsthead
\toprule\noalign{}
~ & source\_name & GB & Bases & Reads \\
\midrule\noalign{}
\endhead
\bottomrule\noalign{}
\endlastfoot
0 & fibroblast & 211.403275 & 553,968,406,500 & 1,846,561,355 \\
1 & pachytene spermatocyte & 274.835160 & 715,656,614,700 &
2,385,522,049 \\
2 & round spermatid & 243.128044 & 655,938,457,200 & 2,186,461,524 \\
3 & sperm & 164.131640 & 428,913,635,400 & 1,429,712,118 \\
4 & spermatogonia & 192.794420 & 518,665,980,300 & 1,728,886,601 \\

\end{longtable}

\section{Handling coolers (Or: preparing
coolers)}\label{handling-coolers-or-preparing-coolers}

\begin{figure}

\centering{

A flowchart showing the pipeline from \texttt{.fastq} to
\texttt{.mcool}. The first 6 steps were done with a Probably BioRender
or Inkscape.

}

\caption{\label{fig-flowchart-handling-coolers}}

\end{figure}%

\subsection{\texorpdfstring{The \emph{gwf} workflow
targets}{The gwf workflow targets}}\label{the-gwf-workflow-targets}

A \emph{gwf} workflow was created to handle the first part of the data
processing, and each accesion number (read pair, mate pair) from the
Hi-C sequencing was processed in parallel, so their execution was
independen on the other samples.

\subsubsection{Downloading the reads}\label{downloading-the-reads}

The reads were downloaded from NCBI SRA portal {[}ref{]} directly to GDK
using \texttt{sra-downloader} {[}ref{]} through docker {[}ref{]} as
\texttt{.fastq.gz} files.

\subsubsection{Handling the reference}\label{handling-the-reference}

The latest reference genome for rhesus macaque (\emph{macaca mulata}),
\emph{rheMac10} (or \emph{Mmul\_10}, UCSC or NCBI naming conventions,
respectively) was downloaded to GDK from UCSC web servers with
\texttt{wget} {[}ref{]}. To use \texttt{bwa} (Burrow Wheeler's Aligner)
{[}ref{]} for mapping, rheMac10 needs to be indexed with both
\texttt{bwa\ index} with the \texttt{-\/-bwtsw} option and
\texttt{samtools\ faidx}, which results in six indexing files for
\texttt{bwa\ mem} to use.

\subsubsection{Mapping paired-end reads}\label{mapping-paired-end-reads}

\subsubsection{Pair and sort the reads}\label{pair-and-sort-the-reads}

\subsubsection{Filter (deduplicate)
pairs}\label{filter-deduplicate-pairs}

\subsubsection{Create interaction matrices
(coolers)}\label{create-interaction-matrices-coolers}

\subsubsection{Pooling samples (Merging
coolers)}\label{pooling-samples-merging-coolers}

The strategy to get the best signal was by pooling the interaction
matrices. \texttt{cooler\ merge} was used to merge all samples in each
sub-folder (cell type) to just one interaction matrix for each cell
type. The reason for that is that we choose to trust Wang et al. (2019)
when they say that compartments are highly reproducible between
replicates, and by merging all replicates, we will have a more robust
signal.

\subsubsection{Create multi-resolution coolers
(zoomify)}\label{create-multi-resolution-coolers-zoomify}

\subsubsection{Matrix balancing (Iterative
correction)}\label{matrix-balancing-iterative-correction}

\subsubsection{Eigendecomposition}\label{eigendecomposition}

\newpage{}

\chapter{Results}\label{results}

Here are the glorious results

\newpage{}

\chapter{Discussion}\label{discussion}

Here is the discussion

\newpage{}

\chapter*{Bibliography}\label{bibliography}
\addcontentsline{toc}{chapter}{Bibliography}

\begingroup
\raggedright

\phantomsection\label{refs}
\begin{CSLReferences}{1}{0}
\bibitem[\citeproctext]{ref-jaenike_sex_2001}
Jaenike, John. 2001. {``Sex {Chromosome} {Meiotic} {Drive}.''}
\emph{Annual Review of Ecology and Systematics} 32 (1): 25--49.
\url{https://doi.org/10.1146/annurev.ecolsys.32.081501.113958}.

\bibitem[\citeproctext]{ref-lieberman_aiden_comprehensive_2009}
Lieberman-Aiden, Erez, Nynke L. Van Berkum, Louise Williams, Maxim
Imakaev, Tobias Ragoczy, Agnes Telling, Ido Amit, et al. 2009.
{``Comprehensive {Mapping} of {Long}-{Range} {Interactions} {Reveals}
{Folding} {Principles} of the {Human} {Genome}.''} \emph{Science} 326
(5950): 289--93. \url{https://doi.org/10.1126/science.1181369}.

\bibitem[\citeproctext]{ref-munch_group_2024}
Munch, Kasper. 2024. {``Munch-Group.''}
\url{https://munch-group.org/research.html}.

\bibitem[\citeproctext]{ref-open2c}
Open Chromosome Collective. 2024. {``Open {Chromosome} {Collective}
({Open2C}).''} Resource. \emph{Open2C}. \url{https://open2c.github.io/}.

\bibitem[\citeproctext]{ref-skov_extraordinary_2023}
Skov, Laurits, Moisès Coll Macià, Elise Anne Lucotte, Maria Izabel Alves
Cavassim, David Castellano, Mikkel Heide Schierup, and Kasper Munch.
2023. {``Extraordinary Selection on the Human {X} Chromosome Associated
with Archaic Admixture.''} \emph{Cell Genomics} 3 (3): 100274.
\url{https://doi.org/10.1016/j.xgen.2023.100274}.

\bibitem[\citeproctext]{ref-sorensen_genome_wide_2023}
Sørensen, Erik F., R. Alan Harris, Liye Zhang, Muthuswamy Raveendran,
Lukas F. K. Kuderna, Jerilyn A. Walker, Jessica M. Storer, et al. 2023.
{``Genome-Wide Coancestry Reveals Details of Ancient and Recent
Male-Driven Reticulation in Baboons.''} \emph{Science} 380 (6648):
eabn8153. \url{https://doi.org/10.1126/science.abn8153}.

\bibitem[\citeproctext]{ref-wang_reprogramming_2019}
Wang, Yao, Hanben Wang, Yu Zhang, Zhenhai Du, Wei Si, Suixing Fan,
Dongdong Qin, et al. 2019. {``Reprogramming of {Meiotic} {Chromatin}
{Architecture} During {Spermatogenesis}.''} \emph{Molecular Cell} 73
(3): 547--561.e6. \url{https://doi.org/10.1016/j.molcel.2018.11.019}.

\end{CSLReferences}

\endgroup


\backmatter


\end{document}
